% Original applied explanation; see Chapter04/SOURCES.md.
\subsubsection{একই লেখা থেকে ব্রাউজার গঠন বোঝে কীভাবে?}
একটি নোটিশে কলেজের নাম, বিজ্ঞপ্তির শিরোনাম ও শেষ তারিখ থাকতে পারে। শুধু কয়েকটি শব্দ বড় করলেই তাদের ভূমিকা নির্দিষ্ট হয় না। \texttt{<h1>} দিয়ে প্রধান শিরোনাম এবং \texttt{<p>} দিয়ে অনুচ্ছেদ চিহ্নিত করলে ব্রাউজার লেখার গঠন বুঝতে পারে।

HTML পড়ার সময় ব্রাউজার এলিমেন্ট ও তাদের ভেতরের লেখা নিয়ে \textbf{Document Object Model (DOM)} নামে নথির গঠন তৈরি করে। সেই গঠন ও প্রযোজ্য রীতি দিয়ে পর্দায় ফল দেখায়। তাই HTML শেখার উদ্দেশ্য কেবল ট্যাগের নাম মুখস্থ করা নয়; কোন তথ্য কোন অংশে থাকবে, সেটি প্রকাশ করতে শেখা।
